\begin{abstract}
A new annotation scheme, language or domain usually starts with a guideline and little annotated data. We propose \emph{decision orchestration}: the guideline is written as an executable procedure in which rules enumerate candidates and assemble the output, and a decision model answers only closed questions, returning a probability for each option the guideline allows. The output is well-formed by construction, every decision is cheap and has a confidence, and uncertain decisions can be routed to a stronger model. We learn the procedure, not model weights: an LLM reflector reads five annotated sentences at a time, with the gold and predicted annotations and every question asked, and proposes edits to the procedure's rules, options and thresholds; each edit is kept only if it measurably improves the score. On verbal multiword expression identification (PARSEME), a procedure learned from 60 annotated sentences reaches 46.3 F in English, 13 points above the best system of the 2018 shared task, which was trained on 3,471 sentences, and 57.8 F in Chinese, 18 points above an LLM annotating directly with ten gold examples; reflection costs under one dollar per language. On Uniform Meaning Representation parsing, where the output is mostly graph structure, closed decisions alone do not beat direct generation, although reflection adapts a Chinese procedure to English (+4.1 F) and routing adds 2.5 F. Decision orchestration pays off where a guideline is a sequence of tests over enumerable candidates.
\end{abstract}
